\uuid{XwR3}
\titre{Pgcd}
\auteur{bodin ; hanani ; mzari}
\organisation{exo7}
\licence{CC BY-NC-SA 4.0}
\datecreate{2020-01-20}
\niveau{L1}
\module{Algèbre}
\chapitre{Polynôme, fraction rationnelle}
\sousChapitre{Division euclidienne}
\difficulte{2}
\theme{Arithmétique des polynômes, pgcd}
\exo7qcmid{lille-1.polynomes.014}

\contenu{
\begin{question}
\qclassification{105.01,105.02}
\qtags{motcle=pgcd}
Quelles sont les affirmations vraies pour des polynômes de $\Rr[X]$ ?
\begin{answers}
    \bad{Le pgcd de $(X-1)^2(X-3)^3(X^2+X+1)^3$ et
    $(X-1)^2(X-2)(X-3)(X^2+X+1)^2$ est $(X-1)^2(X-3)(X^2+X+1)$.}

    \bad{Le ppcm de $(X-1)^2(X-3)^3(X^2+X+1)^3$ et
    $(X-1)^2(X-2)(X-3)(X^2+X+1)^2$ est $(X-1)^2(X-2)(X-3)^3(X^2+X+1)^2$.}

    \good{Le pgcd de $(X-1)^2(X^2-1)^3$ et
    $(X-1)^4(X+1)^5$ est $(X-1)^4(X+1)^3$.}

    \good{Le ppcm de $(X-1)^2(X^2-1)^3$ et
    $(X-1)^4(X+1)^5$ est $(X-1)^5(X+1)^5$.}
\end{answers}
\begin{explanations}
Le pgcd s'obtient en prenant le minimum entre les exposants, le ppcm en prenant le maximum. Attention $X^2-1=(X-1)(X+1)$.
\end{explanations}
\end{question}
}
